\documentclass[aspectratio=169]{beamer}

% --- Configurações de Língua e Tipografia ---
\usepackage[utf8]{inputenc}
\usepackage[T1]{fontenc}
\usepackage[portuguese]{babel}
\usepackage{lmodern}

% --- Temas e Cores ---
\usetheme{Madrid}
\usecolortheme{default}

% --- Pacotes Adicionais ---
\usepackage{graphicx}
\usepackage{booktabs}
\usepackage{listings}
\usepackage{xcolor}

% --- Configuração de Código (Listings) ---
\lstset{
    basicstyle=\ttfamily\small,
    keywordstyle=\color{blue},
    commentstyle=\color{green!60!black},
    stringstyle=\color{orange},
    breaklines=true,
    showstringspaces=false,
    frame=single
}

% --- Informações da Apresentação ---
\title[SafeVault]{SafeVault: Cofre Digital de Documentos}
\subtitle{Segurança e Desenvolvimento Seguro (DESOFS 25/26)}
\author[Grupo 3]{João Loureiro, Gonçalo Barbosa, Miguel Amorim, Diogo Maria}
\institute[ISEP]{Instituto Superior de Engenharia do Porto}
\date{18 de junho de 2026}

% --- Logotipo (Opcional - Comentado para evitar erro se não existir) ---
% \logo{\includegraphics[height=0.8cm]{logo_isep.png}}

\begin{document}

% --- Slide de Título ---
\begin{frame}
    \titlepage
\end{frame}

% --- Sumário ---
\begin{frame}{Sumário}
    \tableofcontents
\end{frame}

\section{Introdução e Objetivos}
\begin{frame}{Visão Geral do SafeVault}
    \begin{columns}
        \column{0.6\textwidth}
        \begin{itemize}
            \item \textbf{O que é?} API REST para gestão segura de documentos sensíveis.
            \item \textbf{Tecnologias:} .NET 9, PostgreSQL, Docker.
            \item \textbf{Foco PBS (Professor):}
            \begin{itemize}
                \item Processo de Desenvolvimento Seguro (S-SDLC).
                \item Resiliência: Reduzir, reparar e detetar danos.
                \item Funções de SO via input do utilizador (Ficheiros).
            \end{itemize}
        \end{itemize}
        
        \column{0.4\textwidth}
        \begin{exampleblock}{Arquitetura}
            Clean Architecture + DDD
            \begin{itemize}
                \item 3 Agregados (User, Vault, Doc)
                \item Auditoria Integrada
            \end{itemize}
        \end{exampleblock}
    \end{columns}
\end{frame}

\section{Desenvolvimento Seguro (S-SDLC)}
\begin{frame}{Ciclo de Vida de Desenvolvimento Seguro}
    \begin{itemize}
        \item \textbf{Automação na Pipeline (GitHub Actions):}
        \begin{itemize}
            \item \textbf{SAST:} CodeQL (análise de vulnerabilidades no código).
            \item \textbf{SCA:} Verificação de pacotes vulneráveis (\texttt{dotnet list package}).
            \item \textbf{Secret Scanning:} Prevenção de fuga de chaves via Gitleaks.
        \end{itemize}
        \item \textbf{Governação de Código:}
        \begin{itemize}
            \item \texttt{CODEOWNERS} e revisões obrigatórias.
            \item Checklists de segurança em cada Pull Request.
        \end{itemize}
    \end{itemize}
\end{frame}

\section{Arquitetura e Segurança}
\begin{frame}{Segurança por Design}
    \begin{columns}
        \column{0.5\textwidth}
        \begin{itemize}
            \item \textbf{Autenticação:} BCrypt (Work Factor 12).
            \item \textbf{Autorização:} RBAC (Admin/Manager/Viewer) com verificação server-side.
            \item \textbf{Integridade:} Hash SHA-256 verificado em cada upload/download.
            \item \textbf{Hardening:} Headers de segurança HTTP e erros sem detalhe interno.
        \end{itemize}
        
        \column{0.5\textwidth}
        \begin{alertblock}{IAST em Runtime}
            Middleware de deteção de padrões suspeitos (SQLi/XSS/Path Traversal) em tempo real.
        \end{alertblock}
    \end{columns}
\end{frame}

\begin{frame}{Resiliência: Reduzir, Reparar e Detetar}
    \begin{itemize}
        \item \textbf{Reduzir:} Rate limiting e lockout de conta após 5 falhas.
        \item \textbf{Reparar (Rollback):} Sistema de versionamento de documentos permite recuperar estados anteriores.
        \item \textbf{Detetar:} 
        \begin{itemize}
            \item Auditoria completa na BD de todas as operações sensíveis.
            \item Logs Serilog rotativos com Correlation ID.
        \end{itemize}
    \end{itemize}
\end{frame}

\section{Testes e Validação}
\begin{frame}{Validação e Resultados}
    \begin{itemize}
        \item \textbf{DAST (OWASP ZAP):} 
        \begin{itemize}
            \item API Scan autenticado via OpenAPI.
            \item \textbf{Resultado:} 0 Falhas / 119 Passagens.
        \end{itemize}
        \item \textbf{Cobertura de Testes:}
        \begin{itemize}
            \item 106 testes unitários e de integração (100\% sucesso).
            \item Testes de regressão focados nas ameaças STRIDE.
        \end{itemize}
        \item \textbf{Conformidade:} Checklist ASVS 5.0 preenchido com evidências.
    \end{itemize}
\end{frame}

\section{Demonstração e Conclusões}
\begin{frame}{Demo: Operações de SO Seguras}
    Foco na manipulação de ficheiros no servidor:
    \begin{enumerate}
        \item \textbf{Validação de Magic Bytes:} Impede upload de ficheiros maliciosos.
        \item \textbf{Canonicalização de Paths:} Previne ataques de Path Traversal.
        \item \textbf{Integridade SHA-256:} Garante que o ficheiro no disco não foi alterado.
        \item \textbf{Auditoria:} Registo instantâneo do acesso ao sistema de ficheiros.
    \end{enumerate}
\end{frame}

\begin{frame}{Conclusões}
    \begin{itemize}
        \item O SafeVault demonstra um compromisso com o desenvolvimento seguro desde a fase de desenho.
        \item A automação (ZAP, CodeQL) permitiu detetar e corrigir vulnerabilidades antes da entrega.
        \item O sistema cumpre os requisitos de resiliência exigidos pelo professor PBS.
    \end{itemize}
    \vspace{0.5cm}
    \centering
    \textbf{Questões?}
\end{frame}

\end{document}
